\section{One regional clipping with an explicit energy separation}

\begin{lemma}[Physical restricted errors for a regional hard truncation]
    \label{lem:alvv_single_region_residuals}
    \lean{TNLean.MPS.AreaLaw.ALVV.exists_single_region_hard_truncation_residuals}
    \notready
    \uses{def:chain_model,thm:disjoint_regional_clipping,
        thm:five_region_original_bond_bounds,thm:consecutive_regional_truncation_low}
    Let $q\ge1$, let $0\le h_i\le I$ be nearest-neighbour terms on an
    $L$-site chain, and take $0\le l\le r\le L$.
    The three intervals ${[0,l)}$, ${[l,r)}$ and ${[r,L)}$ determine their actual
    internal sums $A_0,A_1,A_2$. Let $Y$ be the sum of the original bonds
    crossing at $l$ or $r$, with a coincident bond retained once.
    There are at most two such bonds, and $\sum_i h_i=\sum_j A_j+Y$.

    There exist genuine regional minimum energies $e_0,e_1,e_2$ and one
    common unit vector attaining all three. Set $S=\sum_j e_j$ and
    $H=\sum_i h_i-SI$. Its attained ground energy $E$ belongs to $[0,2]$.
    For $\eta>0$ and $t\ge\eta+7$, put
    \begin{align}
        F   & =\min(A_1,(e_1+t)I)+A_0+A_2+Y-SI,\notag \\
        \xi & =32e^{-(t-\eta)/8}.\notag
    \end{align}
    Then $F\le H$ and the attained ground energy $E'$ of $F$ also belongs
    to $[0,2]$. Write $T_M(c)$ for the closed spectral space of the actual
    centered operator $M$ at energies at most $c$. Both errors satisfy
    \begin{align}
        \|(H-F)|_{T_H(E+\eta)}\| & \le\xi,       &
        \|(H-F)|_{T_F(E+\eta)}\| & \le\xi.\notag
    \end{align}
    These operators use the same scalar shift, and
    $H-F=A_1-\min(A_1,(e_1+t)I)$.
    This is the regional construction of
    \cite[Definition~4.3 and Proposition~4.4(1)]{Arad2016RigorousRG}
    with the displayed sufficient cutoff.
\end{lemma}
\begin{proof}\notready
    \uses{thm:disjoint_regional_clipping,thm:five_region_original_bond_bounds,
        thm:consecutive_regional_truncation_low}
    Apply the five-region decomposition with cuts $l,l,r,r$ and remove the
    two empty regions. The disjoint regional minima give the common test
    vector. The crossing interaction is positive and bounded by $2I$,
    giving both relative ground energies in $[0,2]$.
    Apply the single-region error estimate to the original and middle-clipped
    families with $s=1/8$, $u=\eta+7$ and $\varepsilon=E+\eta$.
    The denominator $u-\varepsilon-2|J|$ is at least one, since $E\le2$
    and $|J|\le2$. The resulting coefficient is at most
    $32e^{-(t-\eta)/8}$. Reindexing the actual orthonormal eigenbasis
    identifies its closed spectral span with the canonically ordered one.
\end{proof}

\begin{lemma}[Ordered energies and relative spectral spaces after one clipping]
    \label{lem:alvv_single_region_spectral_comparison}
    \lean{TNLean.MPS.AreaLaw.ALVV.exists_single_region_hard_truncation_spectral_comparison}
    \notready
    \uses{lem:alvv_single_region_residuals}
    In the setting of Lemma~\ref{lem:alvv_single_region_residuals}, use the
    actual centered operators $H,F$, their ground energies $E,E'$, and
    $\xi=32e^{-(t-\eta)/8}$, with $\eta>0$ and $t\ge\eta+7$.
    Number the eigenvalues of each matrix in decreasing order. Whenever
    $\lambda_i(F)\le E+\eta$,
    \begin{equation}
        \lambda_i(F)\le\lambda_i(H)\le\lambda_i(F)+\xi.\notag
    \end{equation}
    In particular, $E'\le E\le E'+\xi$.
    For any $\delta>0$, set
    \begin{equation}
        \eta'=\eta+\xi+2\xi
        \sqrt{\frac{\lceil\eta/(2\xi)\rceil}{\delta}}.\notag
    \end{equation}
    The enlarged clipped space $T_F(E'+\eta')$ is $\delta$-close to
    $T_H(E+\eta)$, and the enlarged original space $T_H(E+\eta')$ is
    $\delta$-close to $T_F(E'+\eta)$.
    These are the actual centered operators and their own ground-energy
    origins; arbitrary ground degeneracy is allowed. This is the
    sufficient-cutoff consequence of
    \cite[Proposition~4.4 and Lemma~4.5]{Arad2016RigorousRG}, with the
    finite window count corrected as in
    \cite{gap:alvv16_hard_truncation_windows}.
\end{lemma}
\begin{proof}\notready
    \uses{lem:alvv_single_region_residuals}
    The normalized minimum eigenvectors and positive lower certificates
    identify the actual ground energies with the last ordered eigenvalues.
    The derived inequality $F\le H$ gives the lower eigenvalue comparison;
    the clipped low-space restricted error gives its upper comparison.
    Applying the two relative spectral-window comparisons to the actual
    original and clipped residuals gives both closeness statements, with
    the displayed enlargement accounting for the ground shift and the
    ceiling number of windows.
\end{proof}
